
\thispagestyle{empty}

%\textbf{\Huge{Nomenclatura}}
%\normalsize
%\vspace{2cm}

\chapter*{Glosario}

\section*{Terminología}

Aquí deberás poner toda la nomenclatura empleada en la tesis, se recomienda utilizar unidades de sistemas métricos oficiales. Por ejemplo, es común expresar la unidad de tiempo del Sistema Internacional de Unidades (segundos) como seg, cuando la unidad oficial es s. Errores similares son gr para gramos, g o utilizar mayúsculas en Kg, cuando lo correcto es kg. consultar: 

\url{https://www.cenam.mx/siu.aspx} 

\section*{Acrónimos, siglas y abreviaturas}

Los acrónimos son siglas que pueden leerse por sílabas, por ejemplo ONU, INE, CINVESTAV, UNAQ. Las siglas incluyen a los acrónimos y a las siglas como PNPC o PRD que no pueden leerse por sílaba. Las abreviaturas son la versión corta de una palabra y va acompañada de punto y lleva acento en caso de que la palabra completa lo tenga, por ejemplo: México y Méx.


\section*{Notación matemática}

La matemática se apoya en un lenguaje simbólico formal, es decir, que para representar un concepto, una relación u operación se siguen una serie de convenciones propias y reglas. Estos símbolos no son abreviaturas, sino entidades con valor propio y autónomo.

Se sugiere consultar la ISO 80000 (de la Organización Internacional de Normalización) o IEC 80000 (de la Comisión Electrotécnica Internacional),es una guía de estilo para el uso de magnitudes físicas, unidades de medida y fórmulas que las involucran, en documentos de carácter científico o educativo a nivel mundial.

\cleardoublepage 














